\documentstyle[% 


righttag% 


,11pt% 


,amssymb% 


,russian%               remove in Eng. 


,izvru%                 Set izven% in Eng. 


]{amsart} 


 


%       Math definitions 


 


\newcommand{\tts}[1]{} 


 


%\hoffset=-15mm%                  For Eng. 


\setcounter{page}{79} 


\god{2000} 


\nomer{5~(456)} %                        Remove in Eng. 


%\nomer{Vol.\,44, No.\,5}%               Set in Eng. 


%\str{\pageref{begin}--\pageref{end}}%  Set in Eng. 


\udk{514.76} 


 


\author{\it Г.Н.\,Тимофеев} 


% NB:   ^^ remove in Eng. 


\title{Интегрируемая $\pi$-структура с симметричной  


второй ковариантной производной аффинора } 


 


\date{} 


%\pagestyle{myheadings}%         Set in Eng. 


\pagestyle{plain} %              Remove in Eng. 


\begin{document} 


%\thispagestyle{plain}%          Set in Eng. 


\maketitle 


\label{begin} 


 


 


{\bf 1.} Интегрируемая $\pi$-структура с аффинором 


$$ 


a^\sigma_\alpha a^\beta_\sigma=\delta^\beta_\alpha\quad 


(\alpha,\beta,\delta,\dotsc=1,2,\dotsc,N) 


$$ 


определяет [1] композицию $X_n\times X_m$ ($n+m=N$) двух базовых 


многообразий $X_n$ и $X_m$. 


 


В данной статье предполагается, что пространство $W_N$ допускает 


композицию $X_n\times X_m$, для которой вторая ковариантная производная 


аффинора симметрична


\begin{equation} 


\nabla_\mu\nabla_\gamma a^\beta_\alpha= 


\nabla_\gamma\nabla_\mu a^\beta_\alpha. 


\end{equation} 


 


Если $g_{\alpha\beta}$ и $\omega_\gamma$ --- соответственно метрический 


тензор и 


дополнительный вектор пространства $W_N$, то для тензора композиции [1] 


\begin{equation} 


a_{\alpha\beta}=a^\sigma_\alpha g_{\sigma\beta} 


\end{equation} 


вторую ковариантную производную можно представить в виде 


$$ 


\nabla_\mu\nabla_\gamma a_{\alpha\beta}= 


\nabla_\mu\nabla_\gamma a^\sigma_\alpha g_{\sigma\beta}+ 


2\omega_\mu\nabla_\gamma a_{\alpha\beta}+ 


2\omega_\gamma\nabla_\mu a_{\alpha\beta}+ 


2\nabla_\mu\omega_\gamma a_{\alpha\beta}- 


4\omega_\mu\omega_\gamma a_{\alpha\beta}. 


$$ 


Поэтому равенство (1) можно заменить соотношением 


\begin{equation} 


\nabla_{[\mu}\nabla_{\gamma]}a_{\alpha\beta}= 


2\nabla_{[\mu}\omega_{\gamma]}a_{\alpha\beta}. 


\end{equation} 


Следовательно, в пространстве $W_N$ композицию (1) можно определить с 


помощью тензора композиции (2), удовлетворяющего условию (3). 


 


Условие (1) равносильно ([2], с.\,126) соотношению 


\begin{equation} 


-R_{\mu\gamma\alpha}{}^\sigma a^\beta_\sigma+ 


R_{\mu\gamma\alpha}{}^\beta a^\sigma_\alpha=0, 


\end{equation} 


где $R_{\mu\gamma\alpha}{}^\beta=\partial_\mu 


\Gamma_{\gamma\alpha}{}^\beta-\partial_\gamma 


\Gamma_{\mu\alpha}{}^\beta+\Gamma_{\mu\sigma}{}^\beta 


\Gamma_{\gamma\alpha}{}^\sigma-\Gamma_{\gamma\sigma}{}^\beta 


\Gamma_{\mu\alpha}{}^\sigma$ --- тензор кривизны 


пространства. В адаптированной к композиции $X_n\times X_m$ системе 


координат из (4) легко получить равенства 


\begin{gather*} 


R_{\mu\gamma i}{}^{\ov j}=R_{\mu\gamma\ov i}{}^j= 


R_{k l \ov i}{}^{\ov j}=R_{\ov k\,\ov l i}{}^j=0,\\ 


%


\text{(а)}\ R_{i\ov k\,\ov l}{}^{\ov j}= 


R_{i\ov l\,\ov k}{}^{\ov j},\quad 


\text{(б)}\ R_{\ov i k l}{}^j=R_{\ov i l k}{}^j \\ 


%


(i,j,k,\dotsc=1,2,\dotsc,n;\ \; \ov i,\ov j,\ov k,\dotsc= 


n+1,\dotsc,n+m=N). 


\end{gather*} 


Следовательно, отличными от нуля адаптированными компонентами 


тензора кривизны являются только компоненты 


$R_{i j k}{}^l$, $R_{\ov i\,\ov j\,\ov k}{}^{\ov l}$, 


$R_{i\ov j\,\ov k}{}^{\ov l}$, $R_{\ov i j k}{}^l$. 


Тогда в адаптированных координатах условия интегрируемости 


$R_{\mu\gamma\alpha}{}^\sigma g_{\sigma\beta}+ 


R_{\mu\gamma\beta}{}^\sigma g_{\alpha\sigma}= 


-4\nabla_{[\mu}\omega_{\gamma]}g_{\alpha\beta}$ 


основных уравнений 


$\nabla_\gamma g_{\alpha\beta}=2\omega_\gamma g_{\alpha\beta}$ 


превращаются в следующую систему: 


\begin{alignat}{2} 


\text{(а}&) &\quad & R_{k l i}{}^s g_{s j}+ 


R_{k l j}{}^s g_{is}=-4\nabla_{[k}\omega_{l]}g_{i j},\notag\\ 


\text{(б}&) &\quad & R_{\ov k l i}{}^s g_{s j}+ 


R_{\ov k l j}{}^s g_{is}=-4\nabla_{[\ov k}\omega_{l]} 


g_{i j},\notag \\ 


\text{(в}&) &\quad & \nabla_{[\ov k}\omega_{\ov l]} 


g_{i j}=0,\notag \\ 


\text{(г}&) &\quad & R_{k l j}{}^s g_{\ov i s}= 


-4\nabla_{[k}\omega_{l]}g_{\ov i j},\notag \\ 


\text{(д}&) &\quad & R_{\ov k l \ov i}{}^{\ov s}g_{\ov s j}+ 


R_{\ov k l j}{}^s g_{\ov i s}= 


-4\nabla_{[\ov k}\omega_{l]}g_{\ov i j},\\ 


\text{(ж}&) &\quad &R_{\ov k\, \ov l\, \ov i}{}^{\ov s}g_{\ov s j}= 


-4\nabla_{[\ov k}\omega_{\ov l]}g_{\ov i j},\notag\\ 


\text{(з}&) &\quad & \nabla_{[k}\omega_{l]}g_{\ov i\, \ov j}=0,\notag\\ 


\text{(и}&) &\quad & R_{\ov k l\ov i}{}^{\ov s}g_{\ov s\ov j}+ 


R_{\ov k l \ov j}{}^{\ov s}g_{\ov i\ov s}= 


-4\nabla_{[\ov k}\omega_{l]}g_{\ov i\,\ov j},\notag\\ 


\text{(к}&) &\quad & R_{\ov k\,\ov l\,\ov i}{}^{\ov s} 


g_{\ov s\ov j}+R_{\ov k\,\ov l\,\ov j}{}^{\ov s} 


g_{\ov i\ov s}=-4\nabla_{[\ov k}\omega_{\ov l]} 


g_{\ov i\,\ov j}.\notag 


\end{alignat} 


Заметим, что из равенств (5 (б), (д), (и)) соответственно можно 


получить 


\begin{gather} 


R_{\ov k l i}{}^s g_{s j}=-2\nabla_{[\ov k}\omega_{l]} 


g_{i j}-2\nabla_{[\ov k}\omega_{i]}g_{j l}+ 


2\nabla_{[\ov k}\omega_{j]}g_{i l},\\ 


R_{k\ov l\,\ov i}{}^{\ov s}g_{\ov s\ov j}= 


-2\nabla_{[k}\omega_{\ov l]}g_{\ov i\,\ov j}- 


2\nabla_{[k}\omega_{\ov i]}g_{\ov l\,\ov j}+ 


2\nabla_{[k}\omega_{\ov j]}g_{\ov l\,\ov i},\\ 


\nabla_{[\ov k}\omega_{l]}g_{\ov i j}- 


\nabla_{[\ov i}\omega_{l]}g_{\ov k j}- 


\nabla_{[\ov k}\omega_{j]}g_{\ov i l}+ 


\nabla_{[\ov i}\omega_{j]}g_{\ov k j}=0. \notag


\end{gather} 


Применив тождествa Риччи к равенствам (5 (в), (з)), будем иметь 


\begin{gather} 


\nabla_{[k}\omega_{l]}g_{j\ov i}+ 


\nabla_{[l}\omega_{j]}g_{k\ov i}+ 


\nabla_{[j}\omega_{k]}g_{l\ov i}=0,\\ 


\nabla_{[\ov k}\omega_{\ov l]}g_{\ov j i}+ 


\nabla_{[\ov l}\omega_{\ov j]}g_{\ov k i}+ 


\nabla_{[\ov j}\omega_{\ov k]}g_{\ov l i}=0. 


\end{gather} 


 


В системе (5 (г) и (ж)) по значениям первых множителей в левых частях можно 


выделить следующие четыре случая: 


\begin{alignat*}{2} 


& 1.\ \nabla_{[k}\omega_{i]}\ne0, 


\ \; \nabla_{[\ov k}\omega_{\ov i]}\ne0, & \quad 


& 3.\ \nabla_{[k}\omega_{i]}=0, 


\ \; \nabla_{[\ov k}\omega_{\ov i]}\ne0,\\ 


& 2.\ \nabla_{[k}\omega_{i]}\ne0, 


\ \; \nabla_{[\ov k}\omega_{\ov i]}=0, & \quad 


& 4.\ \nabla_{[k}\omega_{i]}=0, 


\ \; \nabla_{[\ov k}\omega_{\ov i]}=0. 


\end{alignat*} 


 


Рассмотрим случай 1. Из 


равенств (г) и (ж) вытекают условия $g_{i j}=g_{\ov i\,\ov j}=0$, 


означающие [1] 


изотропность позиций $P(X_n)$ и $P(X_m)$ композиции. В работе [3] для 


этой 


композиции установлено, что матрица метрического тензора имеет 


блочно-диагональный вид 


\begin{equation} 


(g_{\alpha\beta})= 


\begin{pmatrix} 


0 & g_{i\ov j} \\ g_{\ov i j} & 0 


\end{pmatrix} 


; 


\end{equation} 


пространство четномерно ($N=2n$); число измерений позиций 


одинаково; блоки $(g_{i \ov j})$ и $(g_{\ov i j)})$ невырождены и 


существует тензор $g^{\ov i j}$, 


удовлетворяющий условиям 


\begin{equation} 


g_{\ov i j}g^{j \ov k}=\delta^{\ov k}_{\ov i},\quad 


g_{\ov i j}g^{k\ov i}=\delta^k_j. 


\end{equation} 


Нетрудно показать, что при $n\ne2$ из равенств (8) и (9) вытекают 


свойства $\nabla_{[k}\omega_{l]}=\nabla_{[\ov k}\omega_{\ov l]}=0$, 


противоречащие случаю 1. Следовательно, в 


действительности случай 1 не имеет места. 


 


Поскольку позиции $P(X_n)$ и $P(X_m)$ равноправны, то случаи 2 и 3 


принципиально не отличаются. Поэтому исходная совокупность сводится 


только к случаям 3 и 4. Но в случаях 3 и 4 присутствует одно и тоже 


условие 


$\nabla_{[k}\omega_{l]}=0$. Таким образом, можем сказать: хотя бы одна


составлющая дополнительного вектора является градиентной. 


 


Случай 3. Так как вектор $\omega_i$ не градиентен, то из равенства (5 


(г)) 


вытекает, что $g_{i j}=0$ и матрица $(g_{\alpha\beta})$ метрического 


тензора принимает вид 


\begin{equation} 


(g_{\alpha\beta})= 


\begin{pmatrix} 


0 & g_{i j} \\ g_{\ov j i} & g_{\ov i\,\ov j} 


\end{pmatrix} 


. 


\end{equation} 


Относительно размерностей $m$ и $n$ возможны предположения: 


$m<n$, $m\ge n$. При $m<n$ вычисления показывают, что определитель 


$\det(g_{\alpha\beta})$ равен нулю, что противоречит


условию невырожденности вейлевой метрики. 


Следовательно, имеет место только $m\ge n$. 


При $m=n$ матрица $(g_{i\ov j})$ квадратная невырожденная. Поэтому 


существует тензор $g^{\ov i j}$, удовлетворяющий условиям (11). Тогда 


из 


равенства (9) при $n\ne2$ вытекает градиентность вектора $\omega_{\ov 


i}$ (противоречие). Следовательно, $m\ne n$. 


При $m>n$ матрица $(g_{i\ov j})$ прямоугольная размерности $n\times m$, 


ее ранг 


равен $n$ [3]. Расположим ранговый минор в верхнем углу справа в 


матрице 


$(g_{i\ov j})$. Тогда матрицу (12) метрического тензора можно 


представить в 


виде 


\begin{gather*} 


\begin{pmatrix} 


0 & \vdots & g_{i A} & \vdots & g_{i\ov A}\\ 


\hdotsfor[1.5]{5} \\ 


g_{A i} & \vdots & g_{B A} & \vdots & g_{B \ov A}\\ 


\hdotsfor[1.5]{5} \\ 


g_{\ov A i} & \vdots & g_{\ov A B} & \vdots & g_{\ov B\,\ov A} 


\end{pmatrix} 


\\ 


(A,B,\dotsc=n+1,\dotsc,m; 


\ \; \ov A,\ov B,\dotsc=m+1,\dotsc,m+n). 


\end{gather*} 


 


Анализ соотношения (9) при таком строении метрической матрицы 


приводит к равенству $\nabla_{[\ov i}\omega_{\ov j]}=0$, 


противоречащему исходному. Таким 


образом, и случай 3 не имеет места. 


 


Итак, {\it в вейлевом пространстве, допускающем композицию вида $(3)$, 


компоненты дополнительного вектора градиентны}. 





Если композиция изотропная [3], т.\,е. $g_{i j}=g_{\ov i\,\ov j}=0$, 


то матрица 


метрического тензора имеет вид (10) и блоки $g_{i\ov j}$, $g_{\ov j i}$ 


невырождены. 


Умножив равенство (8) на $g^{\ov i j}$, имеем 


$(n-2)\nabla_{[\ov k}\omega_{l]}=0$. Следовательно, 


{\it пространство, допускающее изотропную композицию вида $(3)$ при 


$n>2$ 


является римановым}. При этом равенство (5 (г)) принимает вид 


$R_{k l j}{}^s g_{\ov i s}=0$, откуда следует 


$R_{k l j}{}^s=0$. Соответственно из (5 (ж)) имеем 


$R_{\ov k\,\ov l\,\ov j}{}^{\ov s}=0$. Таким образом, отличными от нуля 


компонентами тензора 


кривизны являются только 


$R_{l\ov k\,\ov i}{}^{\ov j}$ и $R_{\ov l k i}{}^j$, причем они связаны 


соотношением (5 (д)) 


$$ 


R_{\ov k l j}{}^s g_{\ov i s}- 


R_{l\ov k\,\ov i}{}^{\ov s}g_{\ov s j}=0,\quad 


R_{l\ov k\,\ov i}{}^{\ov j}= 


R_{l\ov i\,\ov k}{}^{\ov j},\quad 


R_{\ov l k i}{}^j=R_{\ov l i k}{}^j. 


$$ 


 


{\bf 2.} В $W_N$ имеет место естественное оснащение ([2], [3]) позиций. 


Внутренняя связность, индуцированная на позициях, есть связность Вейля. 


Адаптированные компоненты поляритета и дополнительного вектора на 


позициях совпадают с соответствующими компонентами основных 


тензоров пространства. Если $g_{i j}$ и $g_{\ov i\,\ov j}$ невырождены, 


то коэффициенты 


$G^k_{i j}$ и $G^{\ov k}_{\ov i\,\ov j}$ внутренних связностей позиций 


выражаются с помощью формул [3] 


$$ 


G^k_{i j}=\Gamma_{i j}{}^k+\Gamma_{i j}{}^{\ov s} 


g_{\ov s p}g^{p k},\quad 


G^{\ov k}_{\ov i\,\ov j}=\Gamma_{\ov i\,\ov j}{}^{\ov k} 


+\Gamma_{\ov i\,\ov j}{}^s g_{s\ov p}g^{\ov p\ov k}. 


$$ 


В дальнейшем предполагается, что адаптированные поляритеты на 


позициях не вырождены. Тогда из соотношений (6), (7) можно определить 


компоненты 


\begin{equation} 


\begin{split} 


\text{(а)}\ R_{\ov k l i}{}^j&= 


-2\nabla_{[\ov k}\omega_{l]}\delta^j_i- 


2\nabla_{[\ov k}\omega_{i]}\delta^j_l+ 


\nabla_{[\ov k}\omega_{s]}g^{s j}g_{i l};\\ 


\text{(б)}\ R_{k\ov l\,\ov i}{}^{\ov j}&= 


-\nabla_{[k}\omega_{\ov l]}\delta^{\ov j}_{\ov i}- 


2\nabla_{[k}\omega_{\ov i]}\delta^{\ov j}_{\ov l}+ 


\nabla_{[k}\omega_{\ov s]}g^{\ov s\ov j}g_{\ov i\,\ov l}, 


\end{split} 


\end{equation} 


где они связаны соотношением (5 (д)) с помощью величин $g_{i\ov j}$. 


Здесь следует 


рассмотреть два случая. 


 


а) $g_{i\ov j}=0$. Это условие характеризует [4] композицию как 


ортогональную (трансверсальные позиции взаимно ортогональны). Из 


соотношения (13) следует, что коэффициенты внутренних связностей 


позиций совпадают с соответствующими адаптированными компонентами 


коэффициентов связности пространства. Геометрии на позициях 


римановы. 


 


б) $g_{i\ov j}\ne0$. В этом случае $n>2$ и $m>2$ и из (5 (д)) следует 


$\nabla_{[l}\omega_{\ov i]}=0$. 


Следовательно, при неортогональности допускаемой композиции 


рассматриваемое пространство риманово. Ясно, что отличными от нуля 


адаптированными компонентами тензора кривизны являются только 


$R_{k l i}{}^j$ и $R_{\ov k\,\ov l\,\ov i}{}^{\ov j}$. 


 


\begin{thebibliography}{9} 


 


\bibitem{1} 


Тимофеев Г.Н. {\em Инвариантные признаки специальных композиций в 


пространствах Вейля} // Изв. вузов. Математика. -- 1976. 


-- \No\,1. -- С.\,87--99. 


\bibitem{2} 


Норден А.П. {\em Пространства аффинной связности}. -- 


М.: Наука, 1976. -- 432 с. 


\bibitem{4} 


Леонтьев Е.К. {\em О чебышевских композициях} // Учен. зап. Казанск.


ун-та. -- 1966. -- Т.\,126. -- \No\,1. -- С.\,23--40. 


\bibitem{3} 


Тимофеев Г.Н. {\em Ортогональные композиции в пространствах 


Вейля} // Изв. вузов. Математика. -- 1976. -- \No\,3. -- С\,.73--85. 


 


\end{thebibliography} 


 


\vskip 2cm 


 


\post{\it Марийский государственный}{\it Поступила} 


\post{\it университет}{08.02.1999} 


 


\label{end} 


\end{document} 


